\section{An exterior extension determined by a fixed collar}
\label{sec:collar}

A nearly flat exterior can carry a Ricci-flat field whose size is much
larger than its Ricci tensor.  Cutting the metric off directly to the
Euclidean metric would lose this distinction.  We construct an extension
which retains the field detected by a fixed collar.  The extension is an
analytic function of that collar metric.  The longer exterior is used
only to estimate the error of this function; neither its coordinates nor
any comparison gauge will have to vary differentiably in time.

\subsection{Norms and the extension statement}
\label{collar:statement}

All lengths in this Section are in collar units.  Let
\(\Gamma\subset O(4)\) be a nontrivial finite group acting freely on
\(S^3\), and put
\[
 C_1=\{1\le |x|\le2\}/\Gamma,\qquad
 \mathcal E=\{|x|\ge1\}/\Gamma.
\]
Calculations are made with equivariant Cartesian components on the
Euclidean cover.  For a fixed-ratio shell \(\Omega_r\), use the norm
\begin{equation}
 \|v\|_{H^s(r)}^2
 =\sum_{j=0}^s r^{2j-4}\int_{\widetilde\Omega_r}|\partial^jv|^2\,dx,
 \qquad
 |v|_{C^s,r}=\sum_{j=0}^s r^j\sup_{\widetilde\Omega_r}|\partial^jv|.
 \label{collar:scaled-norms}
\end{equation}
Here \(\widetilde\Omega_r\) is the lifted shell.  Thus derivatives, but
not the sizes of tensor components, are rescaled.  In particular,
\(\|\mathcal K(v)\|_{H^{s-2}(r)}\) is bounded by
\(Cr^{-2}\|v\|_{H^s(r)}\).  Equivalent fixed enlargements of the shells
will be used without changing notation.  At a boundary they are
one-sided annuli of a fixed positive relative width.  Define
\[
 \|v\|_{H^s_\beta(\mathcal E)}
   =\sup_{r\in\{1,2,4,\ldots\}}r^\beta\|v\|_{H^s(r)}.
\]
On filled Euclidean space the corresponding norm also includes an
ordinary Sobolev norm on the ball of radius two.  These conventions
give constants independent of \(\Gamma\): all estimates are on the
fixed cover and finite-group averaging has norm at most one.
Every relation \(\rho\asymp r\) below ranges over a fixed finite
number of neighboring dyadic shells.  Its comparison width is
chosen once to include the initial shells affected by the collar
correction and is independent of the outer radius \(R\).

Fix \(k=30\) and \(1<q<2\).  Let \(R\) be a large dyadic radius;
\(R\) in this Section does not denote scalar curvature.  Suppose a
metric \(g\), expressed in these coordinates on
\(1\le |x|\le4R\), satisfies
\begin{equation}
 |g-\delta|_{C^{k+2},r}\le\theta_r,
 \qquad \sum_{1\le r\le R}\theta_r\le\sigma,
 \qquad
 \theta_r\ge \|g-\delta\|_{H^k(C_1)}r^{-q}.
 \label{collar:metric-smallness}
\end{equation}
The weights are positive and comparable on neighboring shells, with
fixed comparison constants; the hypotheses include the fixed
enlargements needed below.  A fixed multiplicative constant in the
last inequality can be incorporated into the weights.  Set
\begin{equation}
\begin{split}
 m_1={}&\sup_{C_1}|\Ric(g)|
 +\sum_{1\le r\le R}\sup_{\Omega_r^+}
 \left(\theta_r|\Ric(g)|+
       \sum_{j=1}^{k-2}r^j|\nabla_g^j\Ric(g)|\right)
 +\frac{\theta_R}{R^2},
\end{split}
\label{eq:collar-residual}
\end{equation}
where \(\Omega_r^+\) denotes the specified ample enlargement.  Curvature
norms on the right are intrinsic; since the metrics are uniformly
close to \(\delta\), changing to Cartesian component magnitudes costs
only a fixed factor.

\begin{theorem}[Collar extension]
\label{thm:collar-extension}
There are \(\sigma_*>0\), \(R_*<\infty\), and \(C<\infty\), depending
only on \(q,k\), the fixed shell choices, and the weight comparison
constants, with the following property.  If
\(\sigma\le\sigma_*\) and \(R\ge R_*\), there is a metric \(\bar h\)
on \(\mathcal E\) such that
\begin{equation}
\begin{split}
 \|\bar h-\delta\|_{H^k_q}
   &\le C\|g-\delta\|_{H^k(C_1)},\\
 \|g-\bar h\|_{H^k(C_1)}
    +\|\Ric(\bar h)\|_{H^{k-2}_{q+2}}
   &\le C m_1.
\end{split}
\label{eq:collar-extension}
\end{equation}
The map \(g|_{C_1}\mapsto\bar h-\delta\) is real analytic on a fixed
neighborhood of zero in collar \(H^k\), and its first derivative into
\(H^k_q\) has uniformly bounded operator norm.  It is independent of
\(R\), the weights, and the choice of an exterior realizing the collar
data.  There is also such an analytic extension \(\hat h\), equal to
\(g\) near the inner boundary, for which
\[
 \|\hat h-\delta\|_{H^k_q}\le C\|g-\delta\|_{H^k(C_1)},
 \qquad
 \|\Ric(\hat h)\|_{H^{k-2}_{q+2}}\le Cm_1,
\]
and whose first derivative has the same bound.
\end{theorem}

The last term in Equation~\eqref{eq:collar-residual} cannot simply be
discarded.  A regular vacuum curvature field on a finite annulus is
detected at its outer boundary even when its Ricci tensor vanishes.
The term \(\theta_R/R^2\) retains this information.

\subsection{Two local linear estimates}
\label{collar:local-linear}

The flat compatibility used here is the Saint-Venant relation between
metric strain and linearized curvature. Its place in the Calabi complex
is described in \cite[Section~2.2]{Khavkine2017}. The relation between
compatibility, potentials and Korn's inequality is also developed in
\cite{CiarletCiarlet2005} in dimension three. We give the required
four-dimensional, equivariant estimates with their actual derivative
orders and kernel below.

Write \(Sv\) and \(\mathcal K(v)\) for the flat linearizations of
Ricci and Riemann curvature.  With plus-contraction divergence, set
\(Bv=\partial^iv_{ij}-\tfrac12\partial_j\tr_\delta v\).  Directly
differentiating the coordinate formulas for curvature gives
\begin{equation}
 Sv=-\tfrac12\Delta v+\tfrac12\mathcal L_{(Bv)^\sharp}\delta,
 \qquad BSv=0,
 \qquad
 \Delta\mathcal K(v)=\mathcal T(\partial^2Sv).
 \label{collar:flat-identities}
\end{equation}
Here \(\mathcal T\) is a fixed linear combination of components,
independent of the annulus.  These identities hold distributionally
at the Sobolev orders used below.  Also
\(S\mathcal L_V\delta=\mathcal K(\mathcal L_V\delta)=0\).

\begin{lemma}[Local potential and Korn estimates]
\label{collar:potential}
On a fixed smooth Euclidean annulus, ball, or fixed-width annular
overlap in dimension four, there is a bounded linear extraction
\(v\mapsto P v\in H^{k+1}\) with
\begin{equation}
 \|v-\mathcal L_{Pv}\delta\|_{H^k}
       \le C\|\mathcal K(v)\|_{H^{k-2}},
 \qquad \|Pv\|_{H^{k+1}}\le C\|v\|_{H^k}.
 \label{collar:local-potential}
\end{equation}
For vector fields on these domains,
\begin{equation}
 \inf_{W\in\mathfrak e(4)}\|V-W\|_{H^{k+1}}
           \le C\|\mathcal L_V\delta\|_{H^k},
 \label{collar:korn}
\end{equation}
where \(\mathfrak e(4)=\{x\mapsto Ax+b:A^t=-A\}\).
All choices can be linear, normalized modulo this rigid-motion
space, and equivariant.  After dilation, the right side of the first
inequality in Equation~\eqref{collar:local-potential} is multiplied
by \(r^2\), and vector norms in both inequalities are divided by \(r\).
\end{lemma}

\begin{proof}
We give the potential construction to specify its order and its
kernel.  On any of the stated domains, which are simply connected,
the Neumann gradient projection of a one-form \(\alpha\) writes
\(\alpha=dv+w\), where \(\operatorname{div}w=0\), \(w(\nu)=0\), and
\[
 \|w\|_{H^{s+1}}\le C\|d\alpha\|_{H^s}.
\]
The constant in \(v\) is fixed by its mean.  To recall why there is
no additional term, the div--curl integration-by-parts identity
first gives this inequality with an \(L^2\) remainder; its boundary
term is of order zero.  Compactness removes that remainder on the
orthogonal complement of the kernel.  A kernel element is closed,
hence exact on the stated domain, and its divergence and normal
conditions make its potential a constant Neumann harmonic
function.  The kernel is therefore zero.  Differentiation in smooth
boundary charts, followed by recovery of normal derivatives from
divergence and curl, gives the displayed estimate in every fixed
integer order.  This also makes the projection bounded and linear.

For a symmetric tensor \(v\), form its linearized connection
\[
 C^a_{ij}(v)=\tfrac12
       (\partial_i v_{aj}+\partial_j v_{ai}-\partial_a v_{ij}).
\]
For fixed \(a,j\), its curl in the index \(i\) is a component of
\(\mathcal K(v)\).  Apply the gradient projection with \(s=k-2\)
to obtain
\(C^a_{ij}=\partial_i w^a_j+e^a_{ij}\), with
\(\|e\|_{H^{k-1}}\le C\|\mathcal K(v)\|_{H^{k-2}}\).
Symmetry of \(C^a_{ij}\) in \(i,j\) bounds the curl of \(w^a_j\)
in \(H^{k-1}\) by the same quantity.  A second projection, now
with \(s=k-1\), writes \(w^a_j=\partial_j V^a+e'^a_j\), with
\(\|e'\|_{H^k}\le C\|\mathcal K(v)\|_{H^{k-2}}\).
Consequently \(C(v)-D^2V\) has that bound in \(H^{k-1}\).
Since \(C(\mathcal L_V\delta)=D^2V\) and
\(\partial_i v_{aj}=C^a_{ij}(v)+C^j_{ia}(v)\), the derivative
of \(v-\mathcal L_V\delta\) has that bound in \(H^{k-1}\).
Poincare's inequality controls the tensor modulo a constant symmetric
matrix.  Add the affine vector field with half this matrix as its
derivative.  The result proves the first inequality in
Equation~\eqref{collar:local-potential}.  Each projection has a
bounded linear normalization, and the construction also gives its
second inequality.

For Equation~\eqref{collar:korn}, if
\(e_{ij}=\partial_iV_j+\partial_jV_i\), the identity
\[
 2\partial_i\partial_jV_a
  =\partial_i e_{ja}+\partial_j e_{ia}-\partial_a e_{ij}
\]
controls second and higher derivatives.  Poincare's inequality
then controls \(V\) modulo an affine map, and the symmetric part
of that affine map is controlled by \(e\).  Its remaining ambiguity
is precisely \(\mathfrak e(4)\).  Projecting onto that finite
dimensional space fixes it linearly.  Averaging all constructions
over \(\Gamma\) preserves the inequalities.  Dilation proves the
scaled assertions.
\end{proof}

\begin{lemma}[Newton estimates used below]
\label{collar:newton}
Let \(G(x)=c_4|x|^{-2}\) be the fundamental solution for the chosen
sign of the Euclidean Laplacian.
For \(0<q<2\), its Newton operator is bounded
\[
 \Delta^{-1}:H^{k-2}_{q+2}(\R^4)\longrightarrow H^k_q(\R^4).
\]
For \(1<q<2\), if \(h\in H^{k-3}_{q+3}(\R^4)\) and
\(\int_{\R^4}h=0\) componentwise, then
\begin{equation}
 \|\Delta^{-1}h\|_{H^{k-1}_{q+1}}
             \le C\|h\|_{H^{k-3}_{q+3}}.
 \label{collar:zero-mass-newton}
\end{equation}
Finally, on an annulus or ball of outer radius comparable to \(R\),
if \(\Delta W=H\), boundary strips control \(W\) in scaled
\(H^{k-2}\) by \(b_0\), and
\(c_r=r^2\|H\|_{H^{k-4}(r)}\), then
\begin{equation}
 \sup_r\|W\|_{H^{k-2}(r)}\le C\left(b_0+\sum_rc_r\right).
 \label{collar:green-sum}
\end{equation}
The constant is independent of the number of shells.  A filled ball
has no inner boundary term; its innermost shell is replaced by a
fixed ball.
\end{lemma}

\begin{proof}
These are elementary weighted Newton estimates; compare the weighted
elliptic theory and Newton-kernel argument of
\cite[Theorem~1.7]{Bartnik1986}. We prove them for the shell-supremum
norms used here, including the zero-mass improvement.
Sobolev embedding on unit shells bounds a source in the first
assertion by \(D(1+|y|)^{-q-2}\).  At \(|x|=\rho\ge2\), split
the integral into \(|y|<\rho/2\),
\(\rho/2\le |y|\le2\rho\), and \(|y|>2\rho\).
The first part is bounded by
\(CD\rho^{-2}\int_0^\rho(1+s)^{1-q}\,ds\le CD\rho^{-q}\);
the middle part uses
\(\int_{|z|\le3\rho}|z|^{-2}\,dz\le C\rho^2\);
and the last part is bounded by
\(CD\int_{2\rho}^\infty s^{-q-1}\,ds\).
This proves the required size bound.  The scaled interior estimate
for \(\Delta u=f\),
\[
 \|u\|_{H^s(r)}\le
 C\bigl(\|u\|_{L^\infty(\Omega_r^+)}
          +r^2\|f\|_{H^{s-2}(\Omega_r^+)}\bigr),
\]
gives all the stated Sobolev orders.  The ball near the origin is
handled by the same local estimate.  These local estimates follow
by a cutoff and the Fourier multiplier for \(\Delta\), or by the
usual interior energy estimate and its derivatives.

For the second assertion \(h\) is integrable because \(q>1\), and
zero integral permits the identity
\[
 (G*h)(x)=\int_{\R^4}(G(x-y)-G(x))h(y)\,dy.
\]
On \(|y|<\rho/2\) the kernel difference is at most
\(C\rho^{-3}|y|\), so this part is bounded by
\[
 CD\rho^{-3}\int_0^\rho s^4(1+s)^{-q-3}\,ds
       \le CD\rho^{-q-1}.
\]
On the remaining region, the integral of \(G(x-y)h(y)\) is bounded
as in the first assertion, now by \(CD\rho^{-q-1}\).  The
subtracted term there is bounded by
\(C\rho^{-2}\int_{|y|\ge\rho/2}|h(y)|\,dy
 \le CD\rho^{-q-1}\).
Local elliptic estimates give Equation~\eqref{collar:zero-mass-newton}.

For the last assertion, the Dirichlet Green function in a Euclidean
domain has absolute value at most \(C|x-y|^{-2}\), as follows by
subtracting its boundary harmonic correction from the fundamental
solution and applying the maximum principle.  One shell of forcing
with supremum at most \(Cc_rr^{-2}\) contributes at most \(Cc_r\)
at every point: if \(x\) is near that shell, integrate the kernel
over a ball of radius \(Cr\); if it is far away, use the separation
and the shell volume.  Sobolev embedding supplies the forcing
supremum, since \(k-4>2\).  Componentwise maximum principle and
addition of the shell contributions bound \(\|W\|_\infty\) by
\(C(b_0+\sum c_r)\).  The displayed local elliptic estimate gives
the higher norms on interior shells.  The given boundary strips
supply the remaining norms.  This proof also applies by weak
approximation at the stated Sobolev regularity.
\end{proof}

\subsection{A filled-space generalized inverse and the analytic map}
\label{collar:analytic-map}

Let
\[
 \mathcal X=H^k_q(\mathcal E;\operatorname{Sym}^2),\qquad
 \mathcal Y=H^k(C_1;\operatorname{Sym}^2)
       \mathbin{\times}H^{k-2}_{q+2}(\mathcal E;\operatorname{Sym}^2),
 \qquad Av=(v|_{C_1},Sv).
\]
The pair \((g|_{C_1}-\delta,0)\) need not lie in the range of
\(A\): prescribed collar data need not admit a linearized Ricci-flat
extension.  We first solve a projected nonlinear equation.  The
supplied metric on the finite annulus will then control the residual
left by the projection.

We construct a bounded linear \(L:\mathcal Y\to\mathcal X\) such that
\begin{equation}
 ALA=A.
 \label{collar:generalized-inverse}
\end{equation}
Fix a bounded linear extension \(E\) from collar \(H^k\) to
\(H^k(B_2)\), equal to its input on the entire collar.  Such an
extension is obtained by reflection across the inner smooth
boundary in finitely many boundary charts, with the usual finite
linear combination of reflections matching derivatives through
order \(k\), followed by a partition of unity and an interior
cutoff.  Average it over \(\Gamma\).  Fix also a radial cutoff
\(\chi\) equal to one near radius one and zero before radius two.
For \(d=(v_1,f)\), define
\begin{equation}
 f^\#=
 \begin{cases}
 S(Ev_1),& |x|<1,\\
 \chi Sv_1+(1-\chi)f,&1\le |x|\le2,\\
 f,&|x|>2.
 \end{cases}
 \qquad u=\Delta^{-1}(-2f^\#).
 \label{collar:filled-source}
\end{equation}
The source lies in filled \(H^{k-2}_{q+2}\), with norm at most
\(C\|d\|_{\mathcal Y}\); Lemma~\ref{collar:newton} gives
\(\|u\|_{H^k_q}\le C\|d\|_{\mathcal Y}\).
Apply Lemma~\ref{collar:potential} to \(v_1-u|_{C_1}\), extend
its extracted vector linearly to the exterior, and cut that vector
off outside a fixed annulus.  Write the resulting vector as
\(V_d\), and set
\[
 Ld=u|_{\mathcal E}+\mathcal L_{V_d}\delta.
\]
This operator is bounded.  Its collar residual satisfies
\begin{equation}
 \|v_1-(Ld)|_{C_1}\|_{H^k}
       \le C\|\mathcal K(v_1-u)\|_{H^{k-2}(C_1)},
 \qquad S(Ld)=Su.
 \label{collar:L-residual}
\end{equation}

To prove Equation~\eqref{collar:generalized-inverse}, take \(d=Av\)
and fill \(v\) itself by the same extension \(E(v|_{C_1})\), obtaining
\(V\) on all of \(\R^4\).  The construction gives \(f^\#=SV\)
everywhere, including the interpolation region.  Equation~\eqref{collar:flat-identities}
gives \(Bf^\#=0\).  Thus \(Bu=0\): either commute the constant
coefficient derivative with the Newton integral in distributions,
or observe that \(Bu\) is an entire harmonic field decaying at
infinity and apply the mean-value property.  Consequently
\(Su=f^\#\).  The last identity in Equation~\eqref{collar:flat-identities}
shows that \(\mathcal K(V-u)\) is harmonic on the whole space.
It decays at infinity and hence vanishes by the same mean-value
argument.  The local extraction in Equation~\eqref{collar:L-residual}
is now exact.  We obtain \((Ld)|_{C_1}=v|_{C_1}\) and \(S(Ld)=Sv\),
which proves Equation~\eqref{collar:generalized-inverse}.  Filling
the hole is essential here: a decaying harmonic field on the
punctured exterior need not vanish.

Put \(Q=LA\).  Equation~\eqref{collar:generalized-inverse} implies
\[
 Q^2=Q,\qquad AQ=A.
\]
In particular \(\mathcal X_0=\operatorname{im}Q\) is a closed
Banach subspace.  For small \(h\in\mathcal X_0\), define
\[
 \Phi(h)=QL\bigl(h|_{C_1},\Ric(\delta+h)\bigr).
\]
This map is analytic.  Indeed the inverse metric is given by its
convergent Neumann series, and the coordinate expression for Ricci
is a sum of inverse-metric multiples of \(D^2h\) and products of
\(Dh\).  On every rescaled shell, \(H^{k-2}\) is a multiplication
algebra; inserting weights gives an analytic map
\(H^k_q\to H^{k-2}_{q+2}\).  Its differential at zero is \(S\).
For \(h\in\mathcal X_0\),
\[
 D\Phi(0)h=QLAh=Q^2h=h.
\]
The analytic inverse function theorem on this Banach space
therefore solves
\begin{equation}
 \Phi(\bar h-\delta)=QL(g|_{C_1}-\delta,0)
 \label{collar:analytic-equation}
\end{equation}
in a fixed neighborhood.  It yields the first estimate in
Equation~\eqref{eq:collar-extension} and the asserted bounded
first derivative.  Everything in this definition is fixed by the
collar, \(q\), and the choices of linear operators.

It remains to estimate the error.  Set
\begin{equation}
 d_1=(p_1,f)
    =(g|_{C_1}-\bar h|_{C_1},-\Ric(\bar h)),
 \qquad e=\|d_1\|_{\mathcal Y}.
 \label{collar:error-data}
\end{equation}
Equation~\eqref{collar:analytic-equation} says \(QLd_1=0\).
Applying \(A\), and using \(AQ=A\), gives
\begin{equation}
 ALd_1=0.
 \label{collar:annihilated-data}
\end{equation}
No identity \(LAL=L\) is needed.  We will show that
\(d_1-ALd_1\) is small, and then use
Equation~\eqref{collar:annihilated-data}.

\subsection{A comparison gauge on the long finite annulus}
\label{collar:comparison-gauge}

The analytic extension has already been defined from the collar alone.
To estimate its error, we are free to compare it with the supplied
metric in a separate gauge on the longer annulus. The construction
below must retain constants independent of that annulus's length,
even when its coordinate rotations accumulate.
We first construct a linear projection which removes coordinate
strain without a constant depending on \(R\).  For tensors on
\(1\le |x|\le R\), there are linear operators
\(P_0\) and \(\mathcal D z=\mathcal L_{P_0z}\delta\), with
\(P_0z=0\) on all of \(C_1\), such that
\begin{equation}
 \|z-\mathcal Dz\|_{H^k(r)}
 \le C\left(r^2\max_{\rho\asymp r}
       \|\mathcal K(z)\|_{H^{k-2}(\rho)}
       +\mathbf1_{r\le C_0}\|z\|_{H^k(C_1)}\right).
 \label{collar:long-projection}
\end{equation}
Here and below \(C_0\) is a fixed width, not a growing end length.

Here is the construction.  Take a sequence of roomy dyadic annuli
\(\Omega_j\) and smaller neighboring overlaps with fixed relative
widths.  Lemma~\ref{collar:potential} supplies local vectors \(v_j\)
with
\[
 r_j^{-1}\|v_j\|_{H^{k+1}(r_j)}\le C\|z\|_{H^k(r_j)},
 \quad
 \|z-\mathcal L_{v_j}\delta\|_{H^k(r_j)}\le e_j,
 \quad
 e_j=Cr_j^2\|\mathcal K(z)\|_{H^{k-2}(\Omega_j)}.
\]
On each overlap, Equation~\eqref{collar:korn} writes
\[
 v_{j+1}-v_j=A_jx+w_j,
 \qquad r_j^{-1}\|w_j\|_{H^{k+1}}
                  \le C(e_j+e_{j+1}),
\]
where \(A_j\) is skew and commutes with \(\Gamma\).
There is no equivariant translation: a nonzero invariant vector
would be fixed by each nonidentity group element, contrary to its
free action on the sphere.  The rigid projection can be chosen
orthogonal, so its coefficient also obeys
\[
 |A_j|\le C\bigl(\|z\|_{H^k(\Omega_j)}
                       +\|z\|_{H^k(\Omega_{j+1})}\bigr).
\]
Choose \(B_{j+1}=B_j-A_j\), and write \(W_j=v_j+B_jx\).
Then \(W_{j+1}-W_j=w_j\) on the overlap.  With a radial partition
\(\sum_j\psi_j=1\), let \(W=\sum_j\psi_jW_j\).
The part of its strain coming from partition derivatives is
\[
 \sum_j\bigl(d\psi_j\otimes W_j^\flat
                  +W_j^\flat\otimes d\psi_j\bigr).
\]
On a shell subtract any one active \(W_i\) inside this sum;
\(\sum d\psi_j=0\) cancels it exactly.  Only differences of
neighboring \(W_j\)'s remain, whose norms are controlled by the
nearby \(e_j\)'s.  Thus the residual estimate in
Equation~\eqref{collar:long-projection} holds for \(W\) without its
collar term.  In particular a cumulative rotation \(B_j\) never
appears in a residual bound.

Normalize the global rigid ambiguity by projecting \(W|_{C_1}\)
off the rigid space.  By Equation~\eqref{collar:korn}, its remaining
collar norm is bounded by \(C\|z\|_{H^k(C_1)}\) plus the nearby
curvature residuals.  Subtract a fixed bounded extension of this
entire collar vector, supported in \(r\le C_0\) and equal to it on
all of \(C_1\).  The resulting vector is \(P_0z\); this proves
Equation~\eqref{collar:long-projection}.  If
\(z=\mathcal L_V\delta\) with \(V=0\) on \(C_1\), both terms on
the right vanish.  Thus \(\mathcal Dz=z\) on such strains.
Moreover \(P_0z-V\) is a Euclidean Killing field vanishing on a
whole collar, so it is zero.  Since every \(\mathcal Dz\) is itself
such a strain, \(\mathcal D^2=\mathcal D\).

For the nonlinear step, define
\[
 n(z)=\sup_{r\le R}\theta_r^{-1}\|z\|_{H^k(r)}.
\]
The preceding construction gives more information than
Equation~\eqref{collar:long-projection}.  Namely the local potentials
have size \(Cr\theta_r n(z)\), successive correcting rotations
differ by \(C\theta_r n(z)\), and the first correcting rotation has
size \(C\theta_1n(z)\).  Consequently \(\mathcal D\) is bounded
for \(n\), independently of \(R\).  Radially interpolate the
correcting rotations, with value zero on \(C_1\), to write linearly
\begin{equation}
 P_0z=J(x)x+U(x),\qquad J(x)^t=-J(x),\qquad
 J\gamma=\gamma J\quad(\gamma\in\Gamma),
 \label{collar:rotation-decomposition}
\end{equation}
where \(J=U=0\) on \(C_1\), and
\begin{equation}
 \sup|J|\le C\sigma n(z),\qquad
 \sum_{d=1}^{k+1}\|r^d\partial^dJ\|_{L^\infty(\Omega_r)}
       +r^{-1}\|U\|_{H^{k+1}(r)}\le C\theta_r n(z).
 \label{collar:rotation-bounds}
\end{equation}
To see these bounds explicitly, before the collar correction use
\(J=\sum\psi_jB_j\) and \(U=\sum\psi_jv_j\).
For a derivative of \(J\), subtract a neighboring \(B_i\) using
the differentiated partition identity.  The resulting coefficient
is bounded by a fixed sum of \(|B_j-B_i|\), hence by
\(C\theta_rn(z)\).  The undifferentiated bound is the telescoping
sum \(\sum\theta_rn(z)\).  The collar subtraction and a cutoff of
\(J\) there change these estimates only on finitely many comparable
shells.  Bounds for arbitrarily many fixed derivatives of \(J\)
follow from the smooth fixed partition.  This justifies the
supremum version used in Equation~\eqref{collar:rotation-bounds}.

For \(z\in\operatorname{im}\mathcal D\), define
\begin{equation}
 F_z(x)=e^{J(x)}(x+U(x)).
 \label{collar:orthogonal-map}
\end{equation}
Fix a sufficiently large but fixed ball \(n(z)\le M_0\).
If \(\sigma\) is small in terms of \(M_0\), the map is the identity
on \(C_1\), is equivariant, and is a coordinate map into
\(1\le|x|\le4R\).  Indeed Sobolev embedding and
Equation~\eqref{collar:rotation-bounds} give
\(\|DF_z-\Id\|_\infty\le C M_0\sigma\) and
\(|F_z(x)-x|\le C M_0\sigma|x|\).
Extend it by the identity to the inner ball.  The derivative bound
on this convex ball of radius \(R\) proves injectivity by
integrating along segments.  Its image cannot enter the unit ball
from the exterior, since it already fixes that ball; the upper
radius is less than \(4R\).  This also verifies the boundary and
domain margins needed for composition.

The use of an orthogonal exponential in
Equation~\eqref{collar:orthogonal-map} gives
\begin{equation}
 F_z^*g-\bar h=g-\bar h+z+\mathcal E(z),\qquad
 \mathcal E(0)=0,
 \qquad
 n(\mathcal E(z)-\mathcal E(z'))\le C\sigma n(z-z')
 \label{collar:nonlinear-gauge-error}
\end{equation}
on this fixed ball.  We verify the relative error bound.  Write
\(O=e^J\) and \(\Omega_i=O^{-1}\partial_iO\).  Then
\[
 O^{-1}\partial_iF_z=e_i+\partial_iU+\Omega_i(x+U),
 \qquad \Omega_i=\partial_iJ+O(|J|\,|\partial_iJ|).
\]
The undifferentiated \(O\) cancels exactly from the Euclidean
pullback.  Its linear metric term is the symmetric derivative of
\(U+Jx\): the extra undifferentiated \(J\) in \(D(Jx)\) has zero
symmetric part.  The remaining products have a shell factor
\(\theta_r\) and a second factor bounded by \(C\sigma\).
The same statement for differences follows from linearity of
\(J,U\), the integral formula for a difference of matrix
exponentials, and Sobolev multiplication on rescaled shells.
For the part involving \(g-\delta\), the maps and their
differences have relative \(H^{k+1}\) size at most
\(C\sigma\) and \(C\sigma n(z-z')\).  The mean-value formula
for \((g-\delta)\circ F_z-(g-\delta)\circ F_{z'}\), together
with the \(C^{k+2}\) bound in
Equation~\eqref{collar:metric-smallness}, gives the factor
\(C\theta_r\sigma n(z-z')\).  Differentiating the pullback
formula through order \(k\) produces products with the same
bound; the highest derivative of a map occurs linearly and is
controlled in \(L^2\), while the other factors are controlled by
Sobolev multiplication.  Neighboring-shell comparability covers
the displaced points.  This proves
Equation~\eqref{collar:nonlinear-gauge-error}.  In particular an
undifferentiated accumulated rotation produces no Euclidean metric
error of size \(|J|^2\) on a remote shell.

Equation~\eqref{collar:metric-smallness} and the first estimate for
\(\bar h\) imply \(n(g-\bar h)\le C\).  The equation
\[
 z=-\mathcal D(g-\bar h)-\mathcal D\mathcal E(z)
\]
is a contraction on a sufficiently large fixed \(n\)-ball in
\(\operatorname{im}\mathcal D\), once \(\sigma\) is small.
For its solution set \(F=F_z\) and \(p=F^*g-\bar h\).
Then \(\mathcal Dp=0\), \(p|_{C_1}=p_1\), and
Equation~\eqref{collar:long-projection} gives
\begin{equation}
 \|F^*g-\delta\|_{H^k(r)}\le C\theta_r,
 \qquad
 Y:=\sup_{r\le R}r^{-2}\|p\|_{H^k(r)}\le C(e+X),
 \quad
 X:=\sup_{r\le R}\|\mathcal K(p)\|_{H^{k-2}(r)}.
 \label{collar:gauge-bound}
\end{equation}
This gauge is only a comparison device for the present estimate.
It is not used to define the analytic map in
Equation~\eqref{collar:analytic-equation}.

\subsection{Curvature control and the compatibility residual}
\label{collar:residual-proof}

The comparison gauge controls metric differences by linearized
curvature. We now estimate that curvature and use the filled-space
inverse again to bound the incompatibility of the collar and Ricci
data. Filling the inner ball is what removes an otherwise uncontrolled
inner boundary term.
Put \(T=\Ric(F^*g)\).  The coordinate Ricci formula and
Equation~\eqref{collar:gauge-bound} imply
\begin{equation}
 Sp=T-\Ric(\bar h)+E_2=T+f+E_2,
 \qquad \|E_2\|_{H^{k-2}(r)}\le C\theta_rY.
 \label{collar:quadratic-error}
\end{equation}
Indeed, writing \(\Ric(\delta+v)=Sv+N(v)\), its nonlinear
terms satisfy on a shell
\[
 \|N(v)-N(w)\|_{H^{k-2}(r)}
 \le Cr^{-2}(\|v\|_{H^k(r)}+\|w\|_{H^k(r)})
                     \|v-w\|_{H^k(r)}.
\]
Both first factors are \(O(\theta_r)\); the last factor is at
most \(r^2Y\).  On the collar the sharper bound
\(\|E_2\|_{H^{k-2}(C_1)}\le C\sigma e\) holds because
\(p=p_1\) there.

We will repeatedly use the consequence of
Equation~\eqref{eq:collar-residual}
\begin{equation}
 \|T\|_{H^{k-2}(C_1)}
 +\sum_{r\le R}\sum_{j=1}^{k-2}
       \|r^j\partial^jT\|_{H^0(r)}
 \le Cm_1.
 \label{collar:positive-ricci-derivatives}
\end{equation}
Here no sum of unweighted zeroth-order exterior values is being
claimed.  For completeness, covariant derivatives commute with
pullback.  The first derivative of \(F\) is bounded, so the
intrinsic terms \(F^*(\nabla_g^j\Ric(g))\) are bounded by the
suprema in Equation~\eqref{eq:collar-residual} on the enlarged
shells.  To pass from these tensors to Cartesian derivatives,
expand iteratively \(\partial T=\nabla_{F^*g}T+\Gamma(F^*g)*T\).
A term involving no positive covariant derivative of \(T\)
contains a metric derivative and hence a factor \(\theta_r\);
the other terms are bounded by the positive covariant derivatives
already present in the data.  Products through order \(k-2\)
are controlled in scaled \(L^2\) by the \(H^k\) bound for
\(F^*g\).  This proves the shell sum.  On \(C_1\), the
unweighted zeroth-order term is supplied separately in
Equation~\eqref{eq:collar-residual}.  Smooth approximation of
\(F\) on the roomy shells proves the identities first for smooth
maps and then in these finite Sobolev orders; their bounds require
no uncontrolled higher derivatives of \(F\).

Apply the last identity in Equation~\eqref{collar:flat-identities}
to \(p\).  Its forcing satisfies
\begin{equation}
 \sum_{r\le R}r^2
 \|\mathcal T(\partial^2Sp)\|_{H^{k-4}(r)}
             \le C(m_1+e+\sigma Y).
 \label{collar:first-forcing}
\end{equation}
The contribution of \(T\) is controlled by
Equation~\eqref{collar:positive-ricci-derivatives}; the contribution
of \(f\) is at most \(Ce\sum r^{-q-2}\); and that of \(E_2\)
is at most \(CY\sum\theta_r\).  The inner boundary strips have
\(\mathcal K(p)\)-norm at most \(Ce\), because \(p=p_1\) on
the collar.  Outer boundary strips have norm at most
\(C\theta_R/R^2\) by Equation~\eqref{collar:gauge-bound} and
the lower bound for \(\theta_R\) in
Equation~\eqref{collar:metric-smallness}.  Lemma~\ref{collar:newton}
therefore yields
\[
 X\le C(m_1+e+\sigma Y).
\]
Together with Equation~\eqref{collar:gauge-bound}, and then
absorption of \(C\sigma(X+Y)\), this gives
\begin{equation}
 X+Y\le C(m_1+e).
 \label{collar:first-absorption}
\end{equation}
The two derivatives of \(Sp\) are essential: a constant Ricci
component is killed, so it does not accumulate once per shell.

We now run the fixed linear construction \(L\) on the data
\(d_1=(p_1,f)\) from Equation~\eqref{collar:error-data}.  In
particular \(f^\#\) and \(u\) always mean precisely the filled
source and Newton potential in Equation~\eqref{collar:filled-source}.
Let \(D_0=m_1+\sigma e\).  We claim
\begin{equation}
 \|Bf^\#\|_{H^{k-3}_{q+3}(\R^4)}\le CD_0,
 \qquad \int_{\R^4}Bf^\#\,dx=0.
 \label{collar:bianchi-source}
\end{equation}
Inside radius one, \(Bf^\#=BS(Ep_1)=0\).
On \(C_1\), Equation~\eqref{collar:quadratic-error} and
Equation~\eqref{collar:positive-ricci-derivatives} give
\[
 \|Sp_1-f\|_{H^{k-2}(C_1)}
       \le C(m_1+\sigma e)=CD_0.
\]
Thus derivatives of the fixed blending cutoff cost at most \(CD_0\).
On the exterior, the contracted Bianchi identity for \(\bar h\)
gives \(B_{\bar h}f=0\).  Comparing it with flat \(B\) yields
terms of the form
\((\bar h^{-1}-\delta^{-1})\partial f\) and
\(\partial\bar h*f\), including the corresponding trace terms.
Their weighted \(H^{k-3}_{q+3}\) norm is at most \(C\sigma e\)
by Sobolev multiplication.  This proves the first assertion.
For the second, integrate the divergence expression for each
component of \(Bf^\#\) over a ball.  Its boundary flux is at most
\(Ce\rho^3\rho^{-q-2}=Ce\rho^{1-q}\), which tends to zero
because \(q>1\).  Integrability follows from the first assertion.
There is no inner boundary in this computation.

Commutation with the Newton operator and
Equation~\eqref{collar:zero-mass-newton} now give
\[
 Bu=\Delta^{-1}(-2Bf^\#),\qquad
 \|Bu\|_{H^{k-1}_{q+1}}\le CD_0.
\]
If desired, the commutation identity follows by comparing the
two distributional solutions: their difference is entire
harmonic and decays, hence is zero.  By
Equation~\eqref{collar:flat-identities},
\(Su-f^\#=\tfrac12\mathcal L_{(Bu)^\sharp}\delta\).
On the collar \(f^\#-f=\chi(Sp_1-f)\), while outside it
\(f^\#=f\).  Hence
\begin{equation}
 \|Su-f\|_{H^{k-2}_{q+2}(\mathcal E)}\le CD_0.
 \label{collar:ricci-compatibility}
\end{equation}

To control the other component of the residual, fill \(p\) by
\(Ep_1\) in the unit ball, obtaining \(P\) on the full ball of
radius \(R\).  Consider \(W=\mathcal K(P-u)\).  On the filled
inner ball, \(SP-f^\#=0\); on the collar,
\(SP-f^\#=(1-\chi)(Sp_1-f)\).  Therefore
\(S(P-u)\) on this whole fixed ball is bounded in
\(H^{k-2}\) by \(CD_0\).  Farther out it is
\[
 S(P-u)=T+E_2+(f-Su).
\]
The shell sum of the forcing
\(\Delta W=\mathcal T(\partial^2S(P-u))\) is consequently
bounded by
\[
 C\bigl(m_1+\sigma Y+D_0\bigr)
     \le C(m_1+\sigma e),
\]
where Equation~\eqref{collar:first-absorption} was used in the
last step.  The filled ball contributes one ordinary fixed-scale
term to this sum.  At its only boundary, near radius \(R\),
\begin{equation}
 \|W\|_{H^{k-2}(R)}
 \le C\left(\frac{\theta_R}{R^2}+eR^{-q-2}\right)
 \le C(m_1+eR^{-q-2}).
 \label{collar:outer-curvature}
\end{equation}
Here \(\|u\|_{H^k_q}\le Ce\) follows from the bounded linear
construction of \(L\), and \(\mathcal K(p)\) is controlled on
the outer strip by Equation~\eqref{collar:gauge-bound}.
The filled-ball version of Equation~\eqref{collar:green-sum},
followed by the fixed-scale interior estimate around \(C_1\), gives
\begin{equation}
 \|\mathcal K(p_1-u)\|_{H^{k-2}(C_1)}
       \le C\bigl(m_1+(\sigma+R^{-q-2})e\bigr).
 \label{collar:collar-curvature-residual}
\end{equation}
There is no inner boundary term proportional to \(e\); that is
the purpose of using the same inward extension a second time.

Equations~\eqref{collar:L-residual},
\eqref{collar:ricci-compatibility}, and
\eqref{collar:collar-curvature-residual} imply
\[
 \|d_1-ALd_1\|_{\mathcal Y}
       \le C\bigl(m_1+(\sigma+R^{-q-2})e\bigr).
\]
Since \(ALd_1=0\), first choosing \(\sigma_*\) small and then
\(R_*\) large makes the coefficient of \(e\) at most one half.
Thus \(e\le Cm_1\), proving the second estimate in
Equation~\eqref{eq:collar-extension}.  All absorptions have
constants independent of the exterior length.

Finally choose a fixed smooth radial \(\zeta\), equal to one
near radius one and zero near radius two, and define
\[
 \hat h=\bar h+\zeta(g-\bar h)\quad\hbox{on }C_1,
 \qquad \hat h=\bar h\quad\hbox{for }|x|\ge2.
\]
It is a positive metric for the same smallness choice.  The
coordinate Ricci formula and bounded cutoff derivatives give
\[
 \|\Ric(\hat h)\|_{H^{k-2}(C_1)}
 \le C\bigl(\|\Ric(\bar h)\|_{H^{k-2}(C_1)}
                  +\|g-\bar h\|_{H^k(C_1)}\bigr)\le Cm_1.
\]
Outside the collar the estimate is already proved.  The
interpolation is linear in the collar input and \(\bar h\), so
analyticity and the bounded derivative persist.  This completes
the proof of Theorem~\ref{thm:collar-extension}.

\begin{remark}[Scaling and the time-dependent interface]
\label{collar:scaling-interface}
If the physical collar radius is \(b\), write the normalized
metric as \(g_b=b^{-2}\Phi_b^*g\), where \(\Phi_b(x)=bx\) in
the chosen coordinates.  Covariant Ricci components of \(g_b\)
are \(b^2\) times the physical Cartesian components; intrinsic
Ricci magnitudes and their scaled derivatives consequently gain
the same \(b^2\) factor.  The residual term in physical units is
\(b^{-2}\theta_R/R^2=\theta_R/(Rb)^2\).
For a differentiable family on a fixed material collar, the
bounded derivative of the analytic extension gives
\[
 |\partial_t\hat g|_{0,r}\le CB_1(t)r^{-q},\qquad
 |\Ric(\hat g)|_{0,r}\le Cb^{-2}m(t)r^{-q-2},
\]
when the collar velocity is bounded by \(B_1(t)\) in the
required fixed \(H^k\) norm and \(m(t)\) bounds
Equation~\eqref{eq:collar-residual} in collar units.
The physical volume factor is \(b^4\), so the tail pairing
of these two tensors is bounded by
\(Cb^2m(t)B_1(t)\int_1^\infty r^{1-2q}\,dr\), which is finite
precisely for \(q>1\).  No derivative of the comparison map
\(F\), the far exterior coordinates, the weights, or \(R\) occurs.
All constructions use only the finite orders specified above.
For \(k=30\), scaled Sobolev embedding in dimension four gives
\(C^{27}\) metric bounds and \(C^{25}\) Ricci bounds from
\(H^{30}\) and \(H^{28}\), respectively.  These supply the finite
\(C^{20}\) metric and \(C^{10}\) Ricci bounds used in
Section~\ref{sec:static-input}.
\end{remark}
